% =====================================================================================================
% Chapter 6 - Results and Analysis
% Every number comes from evaluation/results/ (stats.txt, tables/*.csv, replay/, ewm/, speech/, videos/,
% content/, coach/), produced by evaluation/benchmarks/. Figures: docs/report/figures/fig6_*.pdf
% =====================================================================================================
\chapter{Results and Analysis}
\label{ch:results}

\section{Evaluation setup}
\label{sec:setup}

\paragraph{What was evaluated.} The evaluation is a \emph{system-level benchmark without learners}
(Section~\ref{sec:evalproc}). It measures what the system does with a question: whether the answer is
correct and grounded, how long it takes, what it costs in energy and data, and how the adaptive engine
chooses policies as the network changes. It does not measure learning gain or learner satisfaction; those
need a user study (Section~\ref{sec:limitations}).

\paragraph{Hardware and software.} All measurements ran on one laptop acting as the edge device: Intel
Core i5-12450HX, 12~GB RAM, NVIDIA GeForce RTX~5050 Laptop GPU (8~GB), Windows~11, Python~3.11.9. The edge
model \texttt{gemma3:1b} and the embedding model ran on the GPU through Ollama; the cloud models were
\texttt{gpt-oss-120b} (answers) and \texttt{gpt-oss-20b} (grounding check and faithfulness judge) on the
Groq free tier.

\paragraph{Data.} Table~\ref{tab:datasets} summarises what was measured.

\begin{table}[htbp]
  \centering
  \caption{Benchmark data.}
  \label{tab:datasets}
  \begin{tabular}{p{0.28\linewidth}p{0.64\linewidth}}
    \toprule
    Item & Content \\
    \midrule
    Keyed questions & 75 post-test multiple-choice items, 15 per topic (plus 18 pre-test items answered
      while the set was being balanced), each answered by P1--P5, RAG without retry and closed-book
      edge; 47 also by cloud RAG and closed-book cloud. In total 741 measured answers. \\
    Learner doubts & 30 learner-style questions (6 per topic) and 10 clearly off-syllabus questions written
      for the benchmark; 4 variants each, every answer judged for faithfulness (160 answers). \\
    Network profiles & Measured: benchmark network, college Wi-Fi, phone hotspot (Table~\ref{tab:networks}).
      Assumed: a poor rural link and two outages. \\
    Trace replay & 20 repetitions $\times$ 44 sessions $\times$ 12 strategies = 10{,}560 scored sessions
      of 5 answers each. \\
    Other & Speech round trip (40 sentences, 2 recognisers), coach-step timings (2{,}184 repetitions),
      answer scaffolding (180 answers), video strategies (13 runs), 298 content items. \\
    \bottomrule
  \end{tabular}
\end{table}

\begin{table}[htbp]
  \centering
  \caption{Network profiles used in the replay. RTT and loss are measured with 50 TCP connections to the
    Groq API; bandwidth with a timed transfer (the campus firewall blocks speedtest.net).}
  \label{tab:networks}
  \begin{tabular}{lrrrrl}
    \toprule
    Profile & RTT (ms) & Jitter (ms) & Loss & Down / up (Mbps) & Source \\
    \midrule
    Benchmark network (base) & 28.8 & 30.0 & 0\% & -- & measured \\
    Good (college Wi-Fi) & 35.1 & 94.4 & 0\% & 8.7 / 5.8 & measured \\
    Moderate (phone hotspot) & 75.4 & 112.5 & 0\% & 48.2 / 0.8 & measured \\
    Poor (weak rural link) & 600 & 200 & 3\% & 0.5 / 0.25 & assumed \\
    Offline (no route) & -- & -- & -- & -- & cloud call fails at once \\
    Blackhole (no internet) & -- & -- & -- & -- & cloud call waits 60~s, then fails \\
    \bottomrule
  \end{tabular}
\end{table}

\paragraph{Safeguards.} (i) \emph{Paired comparisons}: every strategy is scored on the same questions,
network draws and condition sequence, and the oracle is computed on exactly the same outcomes. (ii)
\emph{Held-out warm start}: warm engines are pre-trained on half of the questions and tested on the other
half. (iii) \emph{No learner data touched}: every run used a copy of the database and its own experience
index. (iv) \emph{Network control}: every answer row records when and on which network it was measured;
31 answers measured after an accidental switch to the hotspot were removed and re-measured on the campus
network. (v) \emph{Fixed seeds} for questions, network draws and exploration. (vi) Scores are computed with
the system's own metric code and fixed weights; Section~\ref{sec:res-ewm} tests other weights.

\section{Evaluation criteria}
\label{sec:criteria}

Table~\ref{tab:criteria} maps the objectives of Section~1.1 to the criteria used.

\begin{table}[htbp]
  \centering
  \caption{Evaluation criteria per objective.}
  \label{tab:criteria}
  \begin{tabular}{p{0.2\linewidth}p{0.45\linewidth}p{0.27\linewidth}}
    \toprule
    Objective & Criterion (unit, interpretation) & Evidence \\
    \midrule
    O1 grounded tutor & accuracy on keyed items (\%); faithfulness: share of claims supported by retrieved
      text (0--1, target $\geq 0.75$); off-topic decline rate (\%) & answer matrix, doubts set \\
    O2 hybrid routing & latency per answer (median and 95th percentile, ms); cloud calls and data per
      answer; energy above idle (mWh) & answer matrix \\
    O3 adaptive selection & session utility $U$ (0--1, higher is better); regret to the oracle (0 is
      perfect); sessions to adapt; prediction error $|\hat U-U|$; answers over 5~s (\%) & trace replay \\
    O4 personalisation & cost of each coach step (ms, model calls); scaffolding overhead (\%); video
      strategy latency, quota, usable videos; content validity & coach, video and content benchmarks \\
    O5 speech & word error rate (\%); recognition and synthesis latency (ms) & speech round trip \\
    Robustness & strategy ranking under EWM weights and $\lambda_1\in[0.5,0.9]$ (Kendall's $\tau$) &
      EWM analysis \\
    \bottomrule
  \end{tabular}
\end{table}

\section{Result}
\label{sec:result}

\subsection{Answer quality: edge versus cloud, with and without retrieval}
\label{sec:res-accuracy}

\begin{figure}[htbp]
  \centering
  \includegraphics[width=0.85\linewidth]{figures/fig6_accuracy_by_retrieval.pdf}
  \caption{Accuracy on the keyed multiple-choice items with 95\% Wilson intervals.}
  \label{fig:accuracy}
\end{figure}

\begin{table}[htbp]
  \centering
  \caption{Accuracy by policy and variant (95\% Wilson interval).}
  \label{tab:accuracy}
  \begin{tabular}{lrrl}
    \toprule
    Variant & $n$ & Accuracy & 95\% CI \\
    \midrule
    P1 edge, RAG $k{=}3$ & 93 & 53.8\% & 43.7--63.5 \\
    P2 edge, CRAG $k{=}5$ & 93 & 51.6\% & 41.6--61.5 \\
    Edge, RAG $k{=}5$ (no retry) & 92 & 52.2\% & 42.1--62.1 \\
    Edge, closed book & 92 & 39.1\% & 29.8--49.3 \\
    P3 cloud, RAG $k{=}3$ & 93 & 100\% & 96.0--100 \\
    P4 cloud, CRAG $k{=}5$ & 92 & 100\% & 96.0--100 \\
    Cloud, closed book & 47 & 97.9\% & 88.9--99.6 \\
    P5 edge with escalation & 92 & 92.4\% & 85.1--96.3 \\
    \bottomrule
  \end{tabular}
\end{table}

The cloud model answered every keyed item correctly; the 1-billion-parameter edge model answered about
half (Figure~\ref{fig:accuracy}, Table~\ref{tab:accuracy}). On the same questions, the cloud was right on
43 items where the edge was wrong and never the reverse (McNemar $p<0.001$). Retrieval helps the edge
model: from 39.1\% closed book to 52.2\% with the lesson context (21 items gained, 9 lost, $p=0.043$),
while the cloud model already knows the material (97.9\% closed book). P5 reaches 92.4\% while escalating
51.1\% of the questions to the cloud after its grounding check (mean check score 0.66): half the cloud
traffic for most of the cloud's accuracy. CRAG's corrective retry never triggered on these items: the
question and its four options always matched the lesson above $\tau=0.65$, so P2 and plain RAG are
identical here; Section~\ref{sec:res-crag} tests CRAG on vaguer questions.

\subsection{Latency and resources per answer}
\label{sec:res-latency}

\begin{figure}[htbp]
  \centering
  \includegraphics[width=0.9\linewidth]{figures/fig6_latency_breakdown.pdf}
  \caption{Median time of one answer by component, on the benchmark network.}
  \label{fig:latency}
\end{figure}

\begin{table}[htbp]
  \centering
  \caption{Latency and resources per answer. Energy is above the idle level (CPU model plus measured GPU
    power); data is what is sent to and received from the cloud.}
  \label{tab:latency}
  \begin{tabular}{lrrrrrr}
    \toprule
    Policy & Median (ms) & p95 (ms) & Retrieval (ms) & Cloud calls & Data (KB) & Energy (mWh) \\
    \midrule
    P1 & 676 & 990 & 45 & 0 & 0 & 6.8 \\
    P2 & 684 & 1093 & 43 & 0 & 0 & 9.2 \\
    P3 & 1342 & 2303 & 45 & 1 & 3.5 & 10.3 \\
    P4 & 1335 & 1938 & 93 & 1 & 5.3 & 6.0 \\
    P5 & 2754 & 3838 & 56 & 1.5 & 8.3 & 18.2 \\
    \bottomrule
  \end{tabular}
\end{table}

On a good network an edge answer takes about 0.7~s and a cloud answer about 1.3~s
(Table~\ref{tab:latency}); both are far inside the 5~s and 10~s targets of NFR-01. Retrieval costs only
40--90~ms; nearly all the time is the language model (Figure~\ref{fig:latency}). P5 is the slowest
(2.75~s median) because its grounding check adds 1.3~s to every edge answer, even when no escalation
follows. Per-answer energy is small (6--18~mWh) and noisy, because the GPU's board power fluctuates
between answers; it is used only at session level in the following sections.

\subsection{Edge versus cloud under each network condition}
\label{sec:res-static}

\begin{table}[htbp]
  \centering
  \caption{Session utility $U$ of each fixed policy under each network profile (mean of 150 sessions per
    cell). The best policy per profile is in bold.}
  \label{tab:static}
  \begin{tabular}{lrrrrr}
    \toprule
    Policy & Good & Moderate & Poor & Offline & Blackhole \\
    \midrule
    P1 edge-light     & 0.747 & 0.755 & 0.751 & 0.744 & \textbf{0.741} \\
    P2 edge-thorough  & 0.748 & 0.752 & 0.755 & \textbf{0.744} & 0.738 \\
    P3 cloud-light    & 0.837 & 0.839 & \textbf{0.820} & 0.705 & 0.582 \\
    P4 cloud-thorough & \textbf{0.838} & \textbf{0.841} & 0.819 & 0.706 & 0.578 \\
    P5 edge-escalate  & 0.808 & 0.808 & 0.774 & 0.706 & 0.578 \\
    \bottomrule
  \end{tabular}
\end{table}

\begin{figure}[htbp]
  \centering
  \includegraphics[width=\linewidth]{figures/fig6_rtt_sweep.pdf}
  \caption{Mean answer latency (left) and session utility (right) as the round-trip time to the cloud
    grows (bandwidth 10/4~Mbps, no loss).}
  \label{fig:rtt}
\end{figure}

Each of P1--P4 is the best fixed choice somewhere (Table~\ref{tab:static}): the cloud policies on good,
moderate and poor networks, the edge policies when the cloud is unreachable. The cost of choosing wrongly
is asymmetric. On a working network the edge loses about 0.09 in utility (accuracy 50\% instead of
100\%); in a blackhole the cloud loses 0.16, because every answer waits 36.7~s on average (up to 61~s)
before the router falls back. The RTT sweep (Figure~\ref{fig:rtt}) shows why the cloud stays preferable
over a wide range: its latency grows from 1.6~s to 9.6~s between 20~ms and 2{,}000~ms of RTT, but its
utility (0.770--0.773 at 2{,}000~ms) remains above the edge's 0.748 because accuracy dominates the score.
P5 falls below the edge at about 1{,}200~ms, because its check adds a second cloud call. The useful
adaptation is therefore not a fine RTT threshold but \emph{recognising when the cloud is unreachable}.

\subsection{Adaptability under a changing network}
\label{sec:res-dynamic}

\begin{figure}[htbp]
  \centering
  \includegraphics[width=0.9\linewidth]{figures/fig6_strategy_utility.pdf}
  \caption{Mean session utility over the dynamic trace with 95\% bootstrap confidence intervals (20
    repetitions). Blue: QAPS; orange: baseline kNN selector; grey: fixed policies, rules and the oracle.}
  \label{fig:strategies}
\end{figure}

\begin{table}[htbp]
  \centering
  \caption{Strategies over the dynamic trace (means per session of 5 answers, 20 repetitions $\times$ 44
    sessions). Regret is the utility gap to the hindsight oracle.}
  \label{tab:dynamic}
  \begin{tabular}{lrrrrrr}
    \toprule
    Strategy & $U$ (95\% CI) & Regret & Accuracy & Latency (s) & $>$5 s & Energy (mWh) \\
    \midrule
    Oracle (hindsight) & 0.818 (0.815--0.820) & 0 & 92.5\% & 2.0 & 4.7\% & 37 \\
    \textbf{QAPS, cold} & \textbf{0.812} (0.809--0.815) & \textbf{0.006} & 90.2\% & 2.1 & 5.6\% & 39 \\
    \textbf{QAPS, warm} & \textbf{0.808} (0.805--0.811) & \textbf{0.009} & 89.0\% & 2.2 & 5.5\% & 42 \\
    P4 cloud-thorough & 0.794 (0.792--0.797) & 0.023 & 90.7\% & 5.4 & 10.9\% & 70 \\
    P3 cloud-light & 0.794 (0.791--0.797) & 0.024 & 90.2\% & 5.4 & 10.6\% & 85 \\
    RTT rule & 0.793 (0.790--0.796) & 0.025 & 76.5\% & 1.2 & 0\% & 45 \\
    Baseline kNN, cold & 0.792 (0.788--0.796) & 0.026 & 83.6\% & 3.4 & 7.4\% & 58 \\
    Baseline kNN, warm & 0.784 (0.781--0.787) & 0.034 & 83.0\% & 4.5 & 9.1\% & 71 \\
    P5 edge-escalate & 0.768 (0.766--0.770) & 0.050 & 84.8\% & 6.8 & 22.1\% & 121 \\
    Random & 0.766 (0.762--0.770) & 0.051 & 71.6\% & 3.8 & 8.9\% & 73 \\
    P2 edge-thorough & 0.741 (0.738--0.744) & 0.077 & 47.9\% & 0.8 & 0\% & 47 \\
    P1 edge-light & 0.735 (0.730--0.739) & 0.083 & 42.9\% & 0.7 & 0\% & 38 \\
    \bottomrule
  \end{tabular}
\end{table}

\begin{figure}[htbp]
  \centering
  \includegraphics[width=\linewidth]{figures/fig6_dynamic_timeline.pdf}
  \caption{Session utility over the trace (mean over repetitions). The shading marks the network
    condition. QAPS follows the oracle through every change; the best fixed policy and the baseline
    selector collapse in the blackhole outage.}
  \label{fig:timeline}
\end{figure}

\paragraph{Overall.} The strategies differ significantly (Friedman $\chi^2=3027$, 880 blocks,
$p<0.001$). QAPS is better than every other strategy in every pairwise comparison (Wilcoxon with Holm
correction, all $p<0.001$): by $+0.017$ to $+0.019$ in utility over the best fixed policies, the RTT rule
and the baseline selector, and by $+0.077$ over the default edge policy P1
(Figure~\ref{fig:strategies}, Table~\ref{tab:dynamic}). Its regret, 0.006, is a quarter of the next best
strategy's (0.023). The baseline kNN selector, which implements the original design, is \emph{not}
significantly better than always using the cloud (P4, $\Delta U=-0.002$, $p=1.0$) or the RTT rule
($p=1.0$).

\begin{table}[htbp]
  \centering
  \caption{Mean session utility per network condition.}
  \label{tab:byprofile}
  \begin{tabular}{lrrrrr}
    \toprule
    Strategy & Good & Moderate & Poor & Offline & Blackhole \\
    \midrule
    Oracle & 0.841 & 0.838 & 0.814 & 0.754 & 0.745 \\
    QAPS, cold & 0.837 & 0.834 & 0.808 & 0.741 & \textbf{0.730} \\
    QAPS, warm & 0.834 & 0.830 & 0.808 & 0.736 & 0.722 \\
    P4 cloud-thorough & 0.839 & 0.836 & 0.811 & 0.698 & 0.566 \\
    RTT rule & 0.838 & 0.835 & 0.729 & 0.737 & 0.720 \\
    Baseline kNN, cold & 0.827 & 0.828 & 0.788 & 0.717 & 0.646 \\
    Baseline kNN, warm & 0.821 & 0.821 & 0.796 & 0.700 & 0.600 \\
    \bottomrule
  \end{tabular}
\end{table}

\paragraph{Where the difference comes from.} Table~\ref{tab:byprofile} shows that each simple strategy
fails in one condition. Always-cloud is near-optimal on working networks but loses 0.18 in the blackhole
outage. The RTT rule protects against outages but abandons the cloud on the poor network, where the cloud
is still better (0.729 instead of 0.811). QAPS (cold) stays within 0.004--0.015 of the oracle in every
condition. The reason is visible in its choices: in offline and blackhole sessions QAPS chose an edge
policy in 100\% of sessions (P1 66--68\%, P2 32--34\%), while the baseline selector still chose a cloud
policy in 46\% (blackhole) and 54\% (offline) of them. Because QAPS never tried the cloud during an
outage, it had no failed cloud calls at all, against 0.27--0.55 per session for the baseline and the
fixed cloud policies, and used 44\% less energy per session than always-cloud (39 versus 70~mWh), close
to the oracle (37~mWh).

\paragraph{Learning and reaction.} After a change of network condition, QAPS chose the condition's best
fixed policy immediately in most cases (median 0 sessions, mean 1.1 for cold, 1.3 for warm) and failed to
find it within the period in only 7--10\% of changes; the baseline needed a median of 1--2 sessions (mean
1.9--2.3) and failed in 21--38\%. QAPS's predictions were also more accurate: mean $|\hat U-U|$ was 0.016
(cold) and 0.036 (warm), against 0.030 and 0.068 for the baseline. The warm-started engines are slightly
worse than the cold ones for both procedures: experience gathered on other questions under the same
network sequence carries over imperfectly, and wrong early beliefs take longer to correct than an empty
memory takes to fill.

\subsection{Corrective retrieval and faithfulness on learner doubts}
\label{sec:res-crag}

\begin{figure}[htbp]
  \centering
  \includegraphics[width=0.95\linewidth]{figures/fig6_doubts_faithfulness.pdf}
  \caption{Left: faithfulness on 30 learner doubts. Right: share of 10 off-topic questions declined.}
  \label{fig:doubts}
\end{figure}

\begin{table}[htbp]
  \centering
  \caption{Learner doubts and off-topic questions. The retry rate equals the share below the CRAG
    threshold because every below-threshold question triggers the retry.}
  \label{tab:doubts}
  \begin{tabular}{llrrrr}
    \toprule
    Set & Variant & Retrieval score & Retried & Faithfulness & Declined \\
    \midrule
    30 doubts & edge RAG & 0.69 & -- & 0.70 & 0\% \\
              & edge CRAG (P2) & 0.69 & 26.7\% & \textbf{0.80} & 0\% \\
              & cloud RAG & 0.69 & -- & 0.99 & 6.7\% \\
              & cloud CRAG (P4) & 0.69 & 26.7\% & 0.98 & 6.7\% \\
    \midrule
    10 off-topic & edge RAG / CRAG & 0.55 & 90\% & 0.43 / 0.50 & \textbf{0\%} \\
                 & cloud RAG / CRAG & 0.55 & 90\% & 1.00 / 1.00 & \textbf{100\%} \\
    \bottomrule
  \end{tabular}
\end{table}

Learner-style doubts are vaguer than test items: 26.7\% of them retrieved below the threshold and
triggered CRAG's corrective retry (Table~\ref{tab:doubts}). On these doubts CRAG raised the faithfulness
of edge answers from 0.70 to 0.80 (Wilcoxon $p=0.027$), above the 0.75 target of objective O1; cloud
answers were already faithful (0.99 and 0.98, no significant difference, $p=0.65$). The off-topic set
exposes a safety gap of the edge model (Figure~\ref{fig:doubts}): retrieval correctly found nothing
relevant for 9 of 10 questions, and the prompt told the model to say so, but the edge model answered all
of them from its own knowledge, whereas the cloud model declined all ten.

\subsection{Cost of the proactive coach}
\label{sec:res-coach}

\begin{table}[htbp]
  \centering
  \caption{Cost of the coach. Top: decision steps (median of repeated timings, no model or network calls).
    Bottom: answer scaffolding on 30 doubts (mean).}
  \label{tab:coach}
  \begin{tabular}{lrrr}
    \toprule
    Coach step & Repetitions & Median & 95th percentile \\
    \midrule
    Plan session goals & 390 & 0.42 ms & 514 ms$^{*}$ \\
    Choose the next step (rules) & 1500 & 0.002 ms & 0.004 ms \\
    Pick two practice items from the pool & 144 & 1.7 ms & 4.0 ms \\
    Score a 15-item post-test & 150 & 0.001 ms & 0.002 ms \\
    \midrule
    Answer style & Edge latency & Edge tokens & Cloud latency \\
    \midrule
    Plain & 1386 ms & 53 & 1762 ms \\
    Step by step & 1733 ms (+25\%) & 112 & 1910 ms (+8\%) \\
    Step by step, hint first & 1860 ms (+34\%) & 122 & 1954 ms (+11\%) \\
    \bottomrule
  \end{tabular}

  \smallskip
  {\footnotesize $^{*}$ when an item has no concept label and one local embedding call is needed.}
\end{table}

The coach's own decisions are effectively free: planning, rule evaluation and scoring take well under a
millisecond and call no model (Table~\ref{tab:coach}). Its cost lies in the actions it triggers. A
scaffolded answer is longer (about twice the tokens on the edge) and 25--34\% slower on the edge, 8--11\%
on the cloud. An accepted ``another video'' suggestion costs a recommendation run (Section~\ref{sec:res-video}).

\subsection{Speech}
\label{sec:res-speech}

\begin{figure}[htbp]
  \centering
  \includegraphics[width=0.7\linewidth]{figures/fig6_speech_wer.pdf}
  \caption{Word error rate and median latency of speech recognition on synthetic speech (Kokoro voices).}
  \label{fig:speech}
\end{figure}

On a round trip from synthetic speech, SraVaani transcribed Hindi with 2.8\% word error and English with
10.7\%, in about 135~ms per question on the GPU; distil-whisper reached 2.6\% on English but took 382~ms
(Figure~\ref{fig:speech}). SraVaani's English errors are mostly scientific terms (\emph{villi}
$\rightarrow$ \emph{villai}, \emph{stomata} $\rightarrow$ \emph{stone ata}, \emph{insulin} $\rightarrow$
\emph{insolin}) and numbers written as words (\emph{9:3:3:1} $\rightarrow$ \emph{nine three three one}).
Recognition is far faster than the 2.6--3.0~s recordings, meeting NFR-02. Speech synthesis is the slow
part: 1.6~s for an English answer (0.6$\times$ real time) but 8.2~s for a Hindi one (2.7$\times$ real
time), which fails the 3~s target for Hindi.

\subsection{Video recommendation}
\label{sec:res-video}

\begin{table}[htbp]
  \centering
  \caption{Recommendation strategies with an empty cache (13 valid runs; each run used 101 YouTube quota
    units).}
  \label{tab:video}
  \begin{tabular}{lrrrrr}
    \toprule
    Strategy & Runs & Time (s) & Candidates & Without English captions & Videos offered \\
    \midrule
    R1 metadata-fast & 5 & 7.1 & 10 & 8.4 (84\%) & 1.6 \\
    R2 balanced & 4 & 7.6 & 15 & 12.8 (85\%) & 2.0 \\
    R3 transcript-deep & 4 & 13.9 & 25 & 14.8 (59\%) & 4.8 \\
    \bottomrule
  \end{tabular}
\end{table}

Most YouTube videos for these chapters are Hindi or Hinglish lectures without an English caption track,
so 84--85\% of the candidates of R1 and R2 had to be dropped and only 1.6--2.0 usable videos remained
(Table~\ref{tab:video}). R3 examines more candidates and offers about three times as many videos, at
twice the time. Making a newly chosen video answerable (fetching, chunking and indexing its transcript)
took 0.4--4.3~s; a video indexed before is reused at once. Seven planned runs could not be completed
because YouTube blocked transcript requests from both test networks after a few hundred requests.

\subsection{Content validity}
\label{sec:res-content}

The five packages contain 298 items over 48 concepts. The blind check agreed with 99.0\% of the answer
keys and 96.0\% of the concept tags. Correct answers are spread evenly over A--D in every topic
(chi-square $p=0.82$--$0.997$), and no item copies more than 9 consecutive words from the textbook. Of the
60 flagged rows, the team decided 32 (29 kept, 3 videos dropped); 28 are still open (3 answer keys, 12
concept tags, 13 videos). No package has yet been formally approved by a reviewer.

\subsection{Sensitivity to the metric weights}
\label{sec:res-ewm}

\begin{table}[htbp]
  \centering
  \caption{Fixed weights versus weights from the Entropy Weight Method (4{,}400 fixed-policy sessions),
    within the measured part of each group.}
  \label{tab:ewm}
  \begin{tabular}{llrr}
    \toprule
    Group & Metric & Fixed & EWM \\
    \midrule
    AQ & retrieval quality / answer quality & 0.57 / 0.43 & 0.42 / 0.58 \\
    NQ & latency / reliability / utilisation / availability & 0.35 / 0.25 / 0.15 / 0.25 & 0.06 / 0.11 / 0.83 / 0.00 \\
    SRE & CPU / memory / energy / network & 0.35 / 0.30 / 0.20 / 0.15 & 0.14 / 0.63 / 0.08 / 0.16 \\
    QoE & AQ / NQ & 0.56 / 0.44 & 0.58 / 0.42 \\
    Utility & QoE / SRE & 0.70 / 0.30 & 0.37 / 0.63 \\
    \bottomrule
  \end{tabular}
\end{table}

EWM assigns weight by variation (Table~\ref{tab:ewm}): device utilisation and memory vary most across
sessions and receive most weight, availability never varied and receives none, and the utility weights
almost reverse (SRE 0.63). When the whole replay, including the engines' own decisions, was repeated with
these weights, QAPS (cold) remained the best strategy after the oracle; QAPS (warm) moved from third to
fifth, behind P4 and P3 by about 0.003; and the baseline selector and RTT rule stayed below QAPS. The
two rankings agree with Kendall's $\tau=0.85$. Varying $\lambda_1$ from 0.5 to 0.9 with the fixed group
weights left QAPS cold and warm in the first two places at every value ($\tau=0.96$--$1.0$ against
$\lambda_1=0.7$).

\section{Comparative Analysis}
\label{sec:comparative}

\paragraph{Against baselines.} Table~\ref{tab:compare} summarises QAPS against the reference
strategies. Relative to the best \emph{fixed} configuration, which is what the problem statement asks to
improve on, QAPS raises utility by 2.3\%, cuts slow answers by half and energy by 44\%, at 0.5 percentage
points lower accuracy. Relative to the system's default policy P1 it gains 10.5\% in utility and 47
points in accuracy. Relative to the original design (baseline selector), it reduces regret by 77\%.

\begin{table}[htbp]
  \centering
  \caption{QAPS (cold) relative to the reference strategies.}
  \label{tab:compare}
  \begin{tabular}{lrrrr}
    \toprule
    Reference & $\Delta U$ & Regret reduction & $\Delta$ answers $>$5~s & $\Delta$ energy \\
    \midrule
    Best fixed policy (P4) & +0.017 (+2.3\%) & 74\% & 10.9\% $\rightarrow$ 5.6\% & $-44\%$ \\
    RTT rule & +0.019 (+2.4\%) & 76\% & 0\% $\rightarrow$ 5.6\% & $-14\%$ \\
    Baseline kNN selector & +0.020 (+2.5\%) & 77\% & 7.4\% $\rightarrow$ 5.6\% & $-33\%$ \\
    Default policy P1 & +0.077 (+10.5\%) & 93\% & 0\% $\rightarrow$ 5.6\% & $+4\%$ \\
    Oracle (upper bound) & $-0.006$ & -- & 4.7\% $\rightarrow$ 5.6\% & $+5\%$ \\
    \bottomrule
  \end{tabular}
\end{table}

\paragraph{Against related work.} The reviewed systems evaluate either the pedagogy or the
infrastructure. AI-agent tutors such as Tebourbi et al.\ (AIAPAL) and Ruffle \& Riley report engagement
and learning outcomes but use one fixed model configuration and do not adapt to the network. Edge
inference work such as Habibi and Er\c{c}etin optimises task completion time and communication overhead
in simulation, without retrieval quality or learners. LAMeTA optimises an intent-aware QoE with
reinforcement learning in a simulated network and reports up to 23.5\% QoE gains over generic DRL; it is
not an educational system. The present evaluation joins the two sides: it measures answer correctness and
grounding together with latency, energy and data under measured and emulated networks, and shows that a
case-based selector without any training reaches 99\% of the hindsight-optimal utility. The numbers are
not directly comparable with those works, which use different tasks, metrics and baselines.

\section{Discussion and Interpretation}
\label{sec:discussion}

\paragraph{O1 -- Grounded tutor.} Retrieval grounding works and CRAG helps where it is meant to: on vague
learner questions it lifted edge faithfulness from 0.70 to 0.80, past the 0.75 target, while on precise
test items it rightly stayed idle. Correctness, however, depends on the model more than on retrieval: the
edge model reaches about 52\% even with the correct passages, the cloud model 100\%. The off-topic result
shows that an instruction to decline is not enough for a 1-billion-parameter model; a rule outside the
model (decline when retrieval stays below the threshold) would close this gap.

\paragraph{O2 -- Hybrid routing.} The edge answers in 0.7~s with no data leaving the device; the cloud in
1.3~s with 3.5--5~KB. P5 shows that a cascade can recover most of the cloud's accuracy (92\%) with half
the cloud calls, but its check makes it the slowest policy, and under this utility it was never the best
choice; it is worth more where cloud calls are expensive or rate-limited.

\paragraph{O3 -- Adaptive selection.} The central claim of the project is supported, with an important
qualification: experience-based policy selection beats every fixed configuration under a changing network
\emph{in the revised form (QAPS)}, not in the form originally designed. The baseline selector's failure is
explained by the measurements: it could not distinguish a poor network from an unreachable cloud,
its shared neighbourhood starved policies of evidence, and its 20\% random exploration kept trying the
cloud during outages. Each QAPS change addresses one of these causes, and the gains are concentrated where
the analysis predicts, in the outage periods. Because utility is dominated by accuracy, the decisive
adaptation is recognising unavailability quickly rather than tuning a latency threshold; this also
explains why a simple RTT rule is a strong but brittle baseline.

\paragraph{O4 -- Personalisation.} The coach's reasoning adds no measurable cost, so proactive
coaching does not compete with the latency budget; its scaffolded answers cost a quarter to a third more
time on the edge, and video suggestions cost a recommendation run. Whether proactive coaching improves
learning, the objective's actual criterion, could not be measured without learners. Video recommendation
is limited mainly by the content available: most suitable lectures lack English captions, which makes the
deeper strategy R3 the only one that reliably offers several usable videos.

\paragraph{O5 -- Speech.} Recognition meets the latency target and is accurate in Hindi; English accuracy
favours distil-whisper, so selecting the recogniser by language would combine the strengths of both.
Hindi speech synthesis is too slow for interactive use on this hardware.

\section{Limitations}
\label{sec:limitations}

\begin{itemize}
  \item \textbf{No learners.} Learning gain, satisfaction and the effect of proactive coaching (EQ and UX
    groups, objective O4) were not measured; benchmark QoE covers answer quality and network/system
    quality only (coverage 0.375).
  \item \textbf{Emulated network dynamics.} All answers and timings are real, but network changes are
    applied by an application-level model; two of five profiles are measured (one session each), the poor
    profile and the outages are assumed, and the model has not yet been validated against live injected
    delay.
  \item \textbf{One device and one case study.} One laptop, one edge model (1B parameters) and one cloud
    provider; five Biology chapters. Results may differ with other hardware, models or subjects.
  \item \textbf{Judges and synthetic inputs.} Faithfulness is judged by an LLM without a human
    agreement check; speech was tested with synthetic voices (a lower bound for real WER); the off-topic
    and Hindi test sentences were written for the benchmark.
  \item \textbf{Content review pending.} Items were drafted with LLM assistance; none is formally
    approved, 28 flags are open, practice and pre-test top-up items are shown without human review, and
    concept-tagging accuracy against two human labellers is not yet available because the labels have not
    been collected.
  \item \textbf{Free-tier limits.} Groq's daily token cap restricted the keyed set to 75 balanced items
    and the closed-book and cloud-RAG controls to 47 items; YouTube blocking left 7 of 20 video runs
    incomplete.
  \item \textbf{Design choices.} The policy set, its levels ($k$, $r$) and the QoE weights are set by
    design; QAPS was developed after observing the baseline's behaviour on the same benchmark, which the
    held-out questions and the weight-sensitivity analysis only partly guard against.
\end{itemize}
